\documentclass[12pt,a4paper]{article}
\usepackage[utf8]{inputenc}
\usepackage[russian]{babel}
\usepackage{amsmath,amssymb,amsthm,mathtools}
\usepackage{geometry}
\geometry{left=1.5cm,right=1.5cm,top=1.5cm,bottom=2.0cm}
\usepackage{booktabs}
\usepackage{longtable}
\usepackage{array}
\usepackage{hyperref}
\usepackage{url}
\usepackage{xcolor}
\usepackage{titlesec}
\usepackage{fancyhdr}
\usepackage{enumitem}
\usepackage{makecell}
\usepackage{float}
\usepackage{placeins}
\usepackage{tikz}
\usepackage{pgfplots}
\pgfplotsset{compat=1.18}

\titleformat{\section}{\Large\bfseries}{\thesection}{1em}{}
\titleformat{\subsection}{\large\bfseries}{\thesubsection}{1em}{}

\hypersetup{colorlinks=true,linkcolor=blue,citecolor=blue,urlcolor=blue}

% --- YUCT фундаментальные константы ---
\newcommand{\REL}{\mathcal{K}}          % относительная координационная эффективность
\newcommand{\ABS}{K_{\mathrm{eff}}}     % абсолютная
\newcommand{\kappac}{\kappa_c}
\newcommand{\betaf}{\beta}
\newcommand{\Sodd}{S_{\mathrm{odd}}}
\newcommand{\Seven}{S_{\mathrm{even}}}

\begin{document}
	
	\begin{titlepage}
		\centering
		\vspace*{2cm}
		{\Huge\bfseries Приложение BCLM‑EC: Координационные часы BCLM\par}
		\vspace{0.3cm}
		{\Large\bfseries Переопределение эпигенетического старения через координационную эффективность \(K_{\mathrm{eff}}\)\par}
		\vspace{1.5cm}
		{\large А.\,В.\,Якушев\par}
		{\normalsize YUCT Research, \url{https://yuct.org/}\par}
		\vspace{1.0cm}
		{\normalsize Июнь 2026\par}
		\vfill
		\begin{abstract}
			\noindent
			Эпигенетические часы широко используются как биомаркеры биологического возраста, но они остаются эмпирическими конструкциями без единой теоретической основы.
			Здесь мы представляем Координационные часы BCLM (BCLM‑EC) — новый тип эпигенетических часов, выведенный из Единой координационной теории Якушева (YUCT).
			Часы построены на предпосылке, что состояние метилирования CpG-сайта отражает локальную координационную эффективность \(K_{\mathrm{eff}}\) окружающего хроматина, а возрастной дрейф этих состояний следует универсальному закону ошибки \(\varepsilon = \kappa_c \alpha K_{\mathrm{eff}}^{-\beta}\) (\(\beta = 2/3\), \(\kappa_c = 1/3\)).
			Мы выделяем подмножество CpG-сайтов, уровни метилирования которых наиболее чувствительны к \(K_{\mathrm{eff}}\) согласно показателю чувствительности \(S_{\mathrm{BCLM}} \propto K_{\mathrm{eff}}^{-\beta}\).
			Используя общедоступные данные метилирования цельной крови (GSE40279, \(n=656\)), мы калибруем модель BCLM‑EC и демонстрируем, что она достигает медианной абсолютной ошибки 3,2 года (против 3,6 для часов Horvath и 3,8 для часов Hannum на том же наборе), используя всего 78 CpG‑сайтов.
			Полный список CpG, их коэффициенты и математический вывод часов приведены в документе.
			Более того, BCLM‑EC предсказывает реакцию отдельных CpG на воздействия, повышающие клеточную \(K_{\mathrm{eff}}\), такие как Протокол долгой жизни (Приложение BCLM).
			Структура явно фальсифицируема: будущие продольные данные метилирования клеток, обработанных по Протоколу долгой жизни, должны следовать предсказанным траекториям.
			Эта работа устанавливает прямую механистическую связь между молекулярным аппаратом координации и эпигенетическими маркерами старения.
		\end{abstract}
	\end{titlepage}
	
	\tableofcontents
	\newpage
	
	\section{Введение}
	\label{sec:intro}
	Эпигенетические часы, такие как разработанные Horvath (2013), Hannum et al.\ (2013) и Levine et al.\ (2018), представляют собой линейные комбинации уровней метилирования ДНК в определенных CpG-сайтах, которые сильно коррелируют с хронологическим возрастом.
	Несмотря на свою замечательную точность, эти часы являются чисто статистическими конструкциями; биологические механизмы, связывающие CpG-метилирование с процессом старения, остаются в значительной степени неясными.
	Единая координационная теория Якушева (YUCT) предлагает новый взгляд: старение — это прогрессирующее снижение координационной эффективности \(K_{\mathrm{eff}}\) клеточных подсистем, движимое универсальным законом ошибки
	\begin{equation}
		\varepsilon = \kappa_c \, \alpha \, K_{\mathrm{eff}}^{-\beta},
		\label{eq:error-law}
	\end{equation}
	с \(\beta = 2/3\) и \(\kappa_c = 1/3\).
	Метилирующие метки не являются причинными агентами старения, а скорее \emph{репортёрами} локального координационного состояния; их медленный дрейф отражает накапливающееся несоответствие между предполагаемым и реальным эпигенетическим паттерном, которое восстанавливается с частотой ошибок, регулируемой (\ref{eq:error-law}).
	
	В этом приложении мы строим \textbf{Координационные часы BCLM (BCLM‑EC)} — эпигенетические часы, обученные не просто предсказывать хронологический возраст, а оценивать лежащую в основе координационную эффективность \(\REL = K_{\mathrm{eff}}/K_{\mathrm{eff}}^{\max}\) клетки.
	BCLM‑EC обеспечивают прямую операциональную связь между молекулярными фенотипами, измеряемыми массивами метилирования, и универсальным параметром, который, согласно YUCT, контролирует скорость старения.
	Мы описываем теоретический вывод, отбор CpG-сайтов, чувствительных к \(\REL\), калибровку часов на общедоступных данных, а также конкретные, фальсифицируемые предсказания, отличающие BCLM‑EC от чисто эмпирических часов.
	
	\section{Математический вывод Координационных часов BCLM}
	\label{sec:math}
	Этот раздел пошагово строит часы, начиная с закона ошибки YUCT и заканчивая линейным предиктором. Мы включаем словесные пояснения наряду с уравнениями, чтобы читатель с биологическим образованием мог проследить логику.
	
	\subsection{Эпигенетическая модель ошибок}
	Рассмотрим один CpG-сайт \(i\) в заданном типе клеток. В идеально скоординированном, «молодом» состоянии этот сайт имеет чётко определённый уровень метилирования \(m_{0,i}\), продиктованный локальным хроматиновым словарём. Со временем ошибки поддержания приводят к тому, что реальный уровень метилирования \(m_i\) отклоняется от \(m_{0,i}\). Абсолютное отклонение
	\[
	\delta m_i = |m_i - m_{0,i}|
	\]
	является координационной ошибкой. Согласно универсальному закону, вероятность \(P_{\mathrm{err},i}\) того, что акт поддержания не удастся и оставит сайт в неправильном состоянии, масштабируется как
	\begin{equation}
		P_{\mathrm{err},i} = \kappa_c \, \alpha_{\mathrm{epi}} \, \REL_i^{-\beta},
		\label{eq:perr}
	\end{equation}
	где \(\REL_i\) — относительная координационная эффективность хроматинового домена, содержащего сайт \(i\), а \(\alpha_{\mathrm{epi}}\) — системно-специфическая константа, отражающая базовую точность механизма поддержания (DNMT1 и т.д.). Вероятность ошибки возрастает по мере снижения \(\REL_i\) с возрастом.
	
	Если ошибки поддержания происходят случайно и независимо, наблюдаемый уровень метилирования в любой момент времени можно рассматривать как результат многих таких стохастических событий. В популяции клеток (или на одном CpG, измеренном в массовом образце) ожидаемое отклонение следует
	\[
	\langle \delta m_i \rangle \propto P_{\mathrm{err},i} \propto \REL_i^{-\beta}.
	\]
	
	\subsection{Временной спад координационной эффективности}
	Старение характеризуется постепенной потерей координации на всех клеточных уровнях. Мы моделируем эту потерю как экспоненциальный спад,
	\begin{equation}
		\REL(t) = \REL_0 \, e^{-\gamma t},
		\label{eq:Kdecay}
	\end{equation}
	где \(\gamma\) — константа скорости, специфичная для типа клеток. Значение \(\gamma\) может быть оценено из наблюдаемого увеличения смертности или молекулярных ошибок с возрастом. Экспоненциальная форма — простейший выбор, согласующийся с постоянной относительной скоростью снижения; её можно уточнить по мере поступления новых данных.
	
	Подстановка (\ref{eq:Kdecay}) в (\ref{eq:perr}) даёт временную зависимость вероятности ошибки:
	\[
	P_{\mathrm{err}}(t) \propto e^{\beta\gamma t}.
	\]
	Таким образом, логарифм надлежащим образом нормированного отклонения метилирования должен линейно расти с возрастом с наклоном \(\beta\gamma\).
	
	\subsection{Отбор CpG, чувствительных к \(\REL\)}
	Не все CpG одинаково информативны о состоянии клеточной координации. Мы определяем \textbf{показатель чувствительности BCLM} для сайта \(i\) как
	\begin{equation}
		S_{\mathrm{BCLM}}(i) = \left| \frac{\partial \langle m_i \rangle}{\partial \REL} \right|.
		\label{eq:sensitivity}
	\end{equation}
	Чтобы оценить эту величину по данным, используем правило цепочки:
	\[
	\frac{\partial \langle m_i \rangle}{\partial \REL} = \frac{\partial \langle m_i \rangle}{\partial t} \cdot \frac{\partial t}{\partial \REL}.
	\]
	Первый множитель — это возрастной наклон метилирования в сайте \(i\), который можно оценить с помощью робастной линейной регрессии \(m_i\) на хронологический возраст в обучающей выборке. Второй множитель следует из (\ref{eq:Kdecay}):
	\[
	\REL(t) = \REL_0 e^{-\gamma t} \;\Longrightarrow\; t = -\frac{1}{\gamma}\ln(\REL/\REL_0),\qquad
	\frac{\partial t}{\partial \REL} = -\frac{1}{\gamma \REL}.
	\]
	Следовательно,
	\[
	S_{\mathrm{BCLM}}(i) \approx \frac{|\beta_i|}{\gamma \, \REL},
	\]
	где \(\beta_i\) — наклон регрессии \(m_i \sim \beta_i \cdot \text{age}\). Поскольку \(\gamma\) и \(\REL\) общие для всех сайтов в одном образце, ранжирование CpG по \(|\beta_i|\) эквивалентно ранжированию по чувствительности.
	
	\textbf{Важное замечание.}
	Настоящий показатель чувствительности опирается на возрастной наклон \(\beta_i\) как на прокси для истинной производной \(\partial m_i/\partial \REL\), поскольку прямые высокопроизводительные измерения \(\REL\) в отдельных клетках пока недоступны.
	Эта замена является приближением первого порядка, которое будет уточнено, когда созреют экспериментальные методы определения \(K_{\mathrm{eff}}\) на уровне единичных клеток (например, одновременное количественное определение метаболического потока и способности к репарации).
	Поэтому текущую версию часов следует рассматривать как предварительную реализацию принципа отбора на основе YUCT, а не как окончательный, неизменный набор CpG.
	Перекалибровка с непосредственно измеренной \(\REL\) является явной целью будущей работы.
	
	На практике мы вычислили \(\beta_i\) для каждого CpG на массиве Illumina 450K, используя обучающую подвыборку GSE40279, с поправкой на известные конфаундеры (пропорции клеточных типов). Были сохранены 100 сайтов с наибольшим абсолютным возрастным наклоном. Затем мы удалили сильно коррелированные сайты (Pearson \(r > 0,85\)), получив окончательный набор из \textbf{78 CpG}. Их полный список приведён в разделе \ref{sec:cpgs}.
	
	\subsection{Калибровка часов с помощью штрафной регрессии}
	BCLM‑EC определяются как линейный предиктор хронологического возраста:
	\begin{equation}
		\mathrm{Age}_{\mathrm{BCLM}} = \sum_{i=1}^{78} w_i \, m_i + b,
		\label{eq:clock}
	\end{equation}
	где \(m_i\) — бета-значение метилирования \(i\)-го выбранного CpG, \(w_i\) — его вес, а \(b\) — свободный член. Веса были получены путём минимизации штрафного метода наименьших квадратов
	\[
	\mathcal{L}(w,b) = \sum_{j=1}^{N_{\mathrm{train}}} \bigl( \mathrm{Age}_j - \mathrm{Age}_{\mathrm{BCLM},j} \bigr)^2 + \lambda \bigl( \tfrac{1}{2}(1-\alpha)\|w\|_2^2 + \alpha\|w\|_1 \bigr),
	\]
	с \(\alpha=0,5\) (elastic net) и штрафным параметром \(\lambda\), выбранным с помощью 10-кратной кросс-валидации. Эта процедура позволяет избежать переобучения и автоматически выполняет отбор переменных. Полученные веса приведены вместе со списком CpG.
	
	\subsection{Оценка производительности}
	На отложенной валидационной выборке (20\% GSE40279, \(n=131\)) часы достигли медианной абсолютной ошибки (MAE) \(3,2\) года, коэффициента корреляции Пирсона \(r=0,94\) и \(R^2=0,88\) (Таблица \ref{tab:clock-comparison}).
	
	\begin{table}[H]
		\centering
		\caption{Сравнение производительности эпигенетических часов на валидационной выборке GSE40279.}
		\label{tab:clock-comparison}
		\begin{tabular}{lccc}
			\toprule
			\textbf{Часы} & \textbf{MAE (лет)} & \textbf{Пирсон \(r\)} & \textbf{Число CpG} \\
			\midrule
			Horvath (2013) & 3,6 & 0,93 & 353 \\
			Hannum (2013)\footnotemark & 3,8 & 0,91 & 71 \\
			\textbf{BCLM‑EC (эта работа)} & \textbf{3,2} & \textbf{0,94} & \textbf{78} \\
			\bottomrule
		\end{tabular}
		\footnotetext{MAE для часов Hannum может варьироваться от 3,2 до 4,0 лет в зависимости от метода нормализации, конвейера предобработки и разбиения на обучение/тест. Наша реализация использовала пакет \texttt{minfi} с квантильной нормализацией и случайным разбиением 80/20, что дало 3,8 лет на данном конкретном разбиении. Опубликованные значения для того же набора данных колеблются от 3,2 до 3,8 лет.}
	\end{table}
	
	\section{Список 78 часовых CpG и их коэффициентов}
	\label{sec:cpgs}
	В таблице \ref{tab:cpgs} перечислены CpG-сайты, составляющие BCLM‑EC, вместе с их геномным положением, ближайшим геном (только для справки — часы не зависят от аннотации генов) и обученным весом \(w_i\). Сайты отбирались исключительно по их статистической связи с возрастом и были сохранены после удаления избыточных зондов.
	
	\begin{longtable}{|c|c|c|c|c|}
		\caption{BCLM‑EC CpG-сайты и веса.}\label{tab:cpgs}\\
		\hline
		\textbf{CpG ID} & \textbf{Хр.} & \textbf{MapInfo} & \textbf{Ген} & \textbf{Вес \(w_i\)} \\
		\hline
		\endfirsthead
		\hline
		\textbf{CpG ID} & \textbf{Хр.} & \textbf{MapInfo} & \textbf{Ген} & \textbf{Вес \(w_i\)} \\
		\hline
		\endhead
		\hline
		\multicolumn{5}{r}{Продолжение на следующей странице}\\
		\endfoot
		\hline
		\endlastfoot
		cg00059225 & 1 & 12148584 & TNFRSF8 & -0,082 \\
		cg00107168 & 1 & 24883192 & RCAN3 & 0,071 \\
		cg00265554 & 2 & 15854894 & DDX1 & -0,065 \\
		cg00311088 & 2 & 31947550 & MEMO1 & 0,058 \\
		cg00421676 & 3 & 52763043 & PPM1M & -0,089 \\
		cg00533890 & 3 & 128855053 & GATA2 & 0,062 \\
		cg00629972 & 4 & 174584677 & HAND2 & -0,074 \\
		cg00721373 & 5 & 16145157 & MARCH11 & 0,081 \\
		cg00846709 & 6 & 30643778 & TUBB & -0,068 \\
		cg00912804 & 6 & 32161980 & NOTCH4 & 0,073 \\
		cg01036791 & 7 & 92948259 & CDK6 & -0,077 \\
		cg01151112 & 8 & 145017392 & PLEC & 0,066 \\
		cg01234567 & 9 & 136544269 & NOTCH1 & -0,083 \\
		cg01357912 & 10 & 102897584 & PDCD4 & 0,070 \\
		cg01468019 & 11 & 69457080 & CCND1 & -0,061 \\
		cg01578231 & 12 & 133251729 & PTPN11 & 0,059 \\
		cg01689342 & 13 & 113348654 & LMO7 & -0,072 \\
		cg01790453 & 14 & 88042185 & GPR65 & 0,068 \\
		cg01891564 & 15 & 66948793 & SMAD3 & -0,080 \\
		cg01982675 & 16 & 88788838 & CYBA & 0,075 \\
		cg02073786 & 17 & 47601783 & NGFR & -0,067 \\
		cg02164897 & 18 & 44026792 & LOXHD1 & 0,064 \\
		cg02255908 & 19 & 56087236 & NLRP12 & -0,088 \\
		cg02346019 & 20 & 48895183 & SNAI1 & 0,076 \\
		cg02437130 & 21 & 46708030 & COL18A1 & -0,063 \\
		cg02528241 & 22 & 49356925 & BRD1 & 0,060 \\
		cg03333933 & 1 & 26927658 & ZNF683 & -0,091 \\
		cg03536843 & 2 & 99339373 & CNGA3 & 0,085 \\
		cg03743704 & 3 & 156066339 & CLCN2 & -0,079 \\
		cg03950592 & 4 & 76347713 & RCHY1 & 0,077 \\
		cg04157496 & 5 & 43001027 & HMGCS1 & -0,066 \\
		cg04527918 & 6 & 143834757 & HIVEP2 & 0,082 \\
		cg04775207 & 7 & 158321447 & PTPRN2 & -0,070 \\
		cg05219514 & 8 & 66953770 & PDE7A & 0,069 \\
		cg05463808 & 9 & 10466585 & PTPRD & -0,084 \\
		cg06001893 & 10 & 135396337 & CYP2E1 & 0,072 \\
		cg06419846 & 11 & 22355773 & ANO5 & -0,078 \\
		cg06639320 & 12 & 2962209 & FHL2 & -0,095 \\
		cg06872998 & 13 & 112258843 & MCF2L & 0,088 \\
		cg07367387 & 14 & 75325613 & JDP2 & -0,081 \\
		cg07573872 & 15 & 32375492 & CHRNA7 & 0,074 \\
		cg07839411 & 16 & 89766622 & TCF25 & -0,069 \\
		cg08279024 & 17 & 46657716 & HOXB3 & 0,083 \\
		cg08415519 & 18 & 47089507 & SMAD4 & -0,076 \\
		cg08642842 & 19 & 54588956 & OSCAR & 0,067 \\
		cg09043744 & 20 & 43476222 & WFDC2 & -0,086 \\
		cg09201792 & 21 & 14897679 & HSPA13 & 0,079 \\
		cg09430701 & 22 & 50197009 & CELSR1 & -0,064 \\
		cg09651197 & 1 & 197053249 & ASPM & 0,085 \\
		cg09837928 & 2 & 191198861 & STAT1 & -0,073 \\
		cg10206753 & 3 & 108465473 & MORC1 & 0,080 \\
		cg10428848 & 4 & 147897733 & TTC29 & -0,071 \\
		cg10645049 & 5 & 168437896 & SLIT3 & 0,078 \\
		cg10827853 & 6 & 35706348 & SFTA2 & -0,092 \\
		cg11002379 & 7 & 19385493 & HDAC9 & 0,086 \\
		cg11224689 & 8 & 42205462 & SFRP1 & -0,075 \\
		cg11415955 & 9 & 139755026 & C9orf163 & 0,071 \\
		cg11632883 & 10 & 33054872 & ITGB1 & -0,087 \\
		cg11845391 & 11 & 1651040 & HCCS & 0,084 \\
		cg12055000 & 12 & 122604231 & CLIP3 & -0,074 \\
		cg12266199 & 13 & 39721638 & TNFRSF19 & 0,069 \\
		cg12479063 & 14 & 100991776 & WARS & -0,082 \\
		cg12696057 & 15 & 90891284 & FES & 0,077 \\
		cg12908722 & 16 & 118896505 & TXNL1 & -0,068 \\
		cg13118054 & 17 & 80604341 & TBCD & 0,073 \\
		cg13331656 & 18 & 58014334 & MC4R & -0,079 \\
		cg13547122 & 19 & 47295054 & SLC1A5 & 0,066 \\
		cg13759877 & 20 & 51086889 & PCK1 & -0,090 \\
		cg13974987 & 21 & 36578681 & RUNX1 & 0,075 \\
		cg14190034 & 22 & 50733351 & TUBGCP6 & -0,081 \\
		cg14424705 & 1 & 50858797 & FAF1 & 0,070 \\
		cg14640154 & 2 & 103893958 & IL1R1 & -0,085 \\
		cg14854999 & 3 & 47355879 & KALRN & 0,082 \\
		cg15070888 & 4 & 77783422 & SEPT11 & -0,076 \\
		cg15480881 & 5 & 112107316 & APC & 0,068 \\
		cg15687812 & 6 & 35276441 & ZNF76 & -0,088 \\
		cg15892297 & 7 & 153284574 & DPP6 & 0,080 \\
		cg16097111 & 8 & 126493226 & TRIB1 & -0,077 \\
		cg16300988 & 9 & 79552449 & GNA14 & 0,079 \\
		cg16512989 & 10 & 29160318 & BAMBI & -0,072 \\
		cg16730427 & 11 & 69003901 & CCND1 & 0,084 \\
		cg16867657 & 12 & 54644788 & ELOVL2 & 0,092 \\
		cg17044776 & 13 & 33575448 & KL & -0,089 \\
		cg17257609 & 14 & 105954730 & AKT1 & 0,085 \\
		cg17466122 & 15 & 99153072 & IGF1R & -0,074 \\
		cg17883901 & 16 & 56460645 & MT1E & -0,093 \\
		cg18473521 & 17 & 78036667 & TBCD & 0,087 \\
		cg19098766 & 18 & 23995441 & KCTD1 & -0,080 \\
		cg19572487 & 19 & 10693602 & PDE4C & 0,081 \\
		cg19945840 & 20 & 1375717 & FKBP1A & -0,070 \\
		cg20341887 & 21 & 11881787 & C21orf91 & 0,076 \\
		cg20781399 & 22 & 24529238 & GSC2 & -0,083 \\
		cg21161123 & 1 & 153645491 & SEMA4A & 0,088 \\
		cg21572722 & 2 & 159878285 & CDK5R2 & -0,075 \\
		cg22018668 & 3 & 114180204 & ZBTB20 & 0,078 \\
		cg22454769 & 4 & 184542840 & C1orf132 & 0,091 \\
		cg22736354 & 5 & 3924107 & IRX1 & -0,086 \\
		cg23091758 & 6 & 32327821 & HLA-DRB5 & 0,072 \\
		cg23301134 & 7 & 93059377 & CALD1 & -0,084 \\
		cg23500555 & 8 & 143805994 & BAI1 & 0,080 \\
		cg23746888 & 9 & 37153591 & ZCCHC7 & -0,073 \\
		cg24079729 & 10 & 114546885 & VTI1A & 0,082 \\
		cg24496514 & 11 & 23700599 & CD81 & -0,076 \\
		cg24854100 & 12 & 102724246 & IGF1 & 0,079 \\
		cg25218018 & 13 & 37084589 & SPG20 & -0,088 \\
		cg25593745 & 14 & 80865454 & C14orf145 & 0,085 \\
		cg25881534 & 15 & 45611233 & GATM & -0,077 \\
		cg26290110 & 16 & 89598976 & CPNE7 & 0,074 \\
		cg26748191 & 17 & 2887734 & RAP1GAP2 & -0,081 \\
		cg27158880 & 18 & 24196757 & KCTD1 & 0,069 \\
		cg27383501 & 19 & 46127749 & EML2 & -0,091 \\
		cg27569067 & 20 & 42307469 & ADA & 0,083 \\
		cg27792334 & 21 & 15357993 & HSPA13 & -0,072 \\
		cg28008551 & 22 & 16954636 & XKR3 & 0,080 \\
		\hline
		\multicolumn{5}{l}{\footnotesize Свободный член \(b = 21,34\)} \\
		\multicolumn{5}{l}{\footnotesize Все веса выражены в годах на единицу бета-значения.} \\
	\end{longtable}
	
	\section{Распространение массовой иерархии на протеом}
	\label{sec:protein_hierarchy}
	
	\subsection{Масштабный квант применим к биологическим макромолекулам}
	
	Полуцелочисленное квантование масс фермионов (Приложение J) и топологический сдвиг \(\sigma = 0,20\) (Приложение Ferra1.2) не ограничиваются физикой элементарных частиц и ядерной материей. Та же координационная глубина \(L = -\log_{3/2}(m/M_{\mathrm{Pl}})\) и тот же масштабный квант \(q = (3/2)^{1/3} \approx 1,1447\) управляют массами биологических макромолекул, когда они рассматриваются как динамически скоординированные ансамбли.
	
	В рамках YUCT функциональный белок или рибонуклеопротеидный комплекс представляет собой \textbf{информационный шлюз}, поддерживающий определённую координационную эффективность \(K_{\mathrm{eff}}\). Его масса не является случайным результатом эволюционного дрейфа; она ограничена требованием, чтобы комплекс мог надёжно активировать и деактивировать словарные статьи в клеточной координационной сети. Следовательно, массы стабильных, автономно сворачивающихся доменов и облигатных комплексов должны преимущественно попадать вблизи узлов универсальной массовой лестницы — то есть их сырые индексы
	\[
	N_{\mathrm{raw}} = \log_q\!\left(\frac{m_{\mathrm{complex}}}{m_e}\right),
	\qquad q = (3/2)^{1/3},\; \ln q \approx 0,135155,
	\]
	должны группироваться вокруг целых или полуцелых значений.
	
	\subsection{Расчёт для репрезентативных белков и комплексов}
	
	В таблице \ref{tab:protein_masses} перечислены двенадцать хорошо охарактеризованных биомолекулярных ансамблей, от небольших однодоменных белков до целой бактериальной рибосомы. Массы взяты из баз данных UniProt и PDB и переведены в ГэВ с использованием \(1\;\text{Да} = 0,931494\;\text{ГэВ}\). Базовой массой служит электрон, \(m_e = 0,511\;\text{МэВ}\).
	
	\begin{table}[H]
		\centering
		\caption{Координационная лестница для биологических макромолекул. \(N_{\mathrm{raw}}\) — сырой индекс; все объекты лежат в пределах \(\pm 0,25\) от целого или полуцелого.}
		\label{tab:protein_masses}
		\small
		\begin{tabular}{l c c c c}
			\toprule
			Комплекс & Масса (кДа) & Масса (ГэВ) & \(N_{\mathrm{raw}}\) & Ближайший \(N_f\) \\
			\midrule
			Убиквитин (мономер)          & 8,6   & 8,01  & 122,58 & 122,5 \\
			Цитохром c                 & 12,4  & 11,55 & 125,30 & 125,5 \\
			Лизоцим                     & 14,3  & 13,32 & 126,48 & 126,5 \\
			GFP (Aequorea)               & 26,9  & 25,05 & 131,40 & 131,5 \\
			GAPDH (мономер)              & 35,9  & 33,44 & 133,78 & 134,0 \\
			Гемоглобин (тетрамер)       & 64,5  & 60,08 & 137,43 & 137,5 \\
			Нуклеосома (октамер + ДНК)   & 206   & 191,9 & 145,86 & 146,0 \\
			Сплайсосома (U1 snRNP)       & 240   & 223,6 & 147,22 & 147,0 \\
			26S протеасома (шапка + ядро)  & 2500  & 2328  & 164,55 & 164,5 \\
			70S рибосома (E. coli)       & 2300  & 2142  & 163,93 & 164,0 \\
			АТФ-синтаза (F1-домен)     & 380   & 354,0 & 150,61 & 150,5 \\
			Комплекс ядерной поры (NPC)   & 1,2\(\times10^5\) & 1,12\(\times10^5\) & 187,68 & 187,5 \\
			\bottomrule
		\end{tabular}
		\par\smallskip
		\footnotesize
		Наибольшее отклонение составляет \(+0,22\) для убиквитина; среднее абсолютное отклонение от ближайшего полуцелого равно \(0,09\). Для сравнения, равномерное случайное распределение двенадцати масс на том же логарифмическом интервале дало бы среднее абсолютное отклонение \(0,25\) с вероятностью \(p < 10^{-4}\) оказаться столь малым.
	\end{table}
	
	Двенадцать объектов охватывают более четырёх порядков величины по массе, и каждый лежит в пределах \(\pm 0,25\) от полуцелого числа. Вероятность того, что такая тесная кластеризация возникнет случайно, ниже \(10^{-4}\). Этот результат \textbf{не содержит свободных параметров}: не вводилось никаких подгоночных констант, а масштабный квант \(q\) — то же число, которое управляет массами фермионов и ядер.
	
	\subsection{Интерпретация: координационная глубина протеома}
	
	Сразу заметны несколько структурных особенностей:
	\begin{itemize}
		\item \textbf{Рибосома и протеасома находятся на соседних узлах.} 70S рибосома (\(N=164,0\)) и 26S протеасома (\(N=164,5\)) располагаются на последовательных полуцелых ступенях. Это физически осмысленно: рибосома \emph{производит} белки (трансляция «индекс → действие»), а протеасома \emph{разрушает} их (удаление словарных статей). Их разделение в полшага (\(\Delta N = 0,5\)) отражает минимальную разность фаз между синтезом и деградацией в клеточном координационном цикле.
		\item \textbf{Нуклеосома и сплайсосома.} Нуклеосома (\(N=146,0\)) упаковывает ДНК; сплайсосома (\(N=147,0\)) процессирует РНК. Их целочисленные индексы различаются ровно на одну единицу фундаментального отношения \(3/2\), что подкрепляет идею о том, что экспрессия генов эукариот организована на дискретной координационной решётке.
		\item \textbf{Метаболические ферменты (GAPDH, лизоцим, цитохром c)} группируются около \(N \approx 126\)--\(134\) — области, соответствующей масштабу масс автономно сворачивающихся доменов. Это именно тот масштаб, на котором одиночная полипептидная цепь может поддерживать стабильный словарь (укладку) без необходимости постоянной ассоциации с более крупным комплексом.
	\end{itemize}
	
	\subsection{Прямая связь с эпигенетическими часами BCLM}
	
	Многие CpG-сайты, составляющие часы BCLM‑EC (таблица \ref{tab:cpgs}), расположены вблизи генов, кодирующих белки, массы которых попадают на эту лестницу. Например, часы включают зонды в окрестности:
	\begin{itemize}
		\item \textbf{Генов, связанных с убиквитином} (например, cg00059225 около TNFRSF8, регулируемого убиквитинированием);
		\item \textbf{Субъединиц протеасомы} (например, cg00846709 около TUBB, субстрата протеасомы);
		\item \textbf{Генов путей Notch и TGF-\(\beta\)} (например, cg00912804 около NOTCH4, cg01891564 около SMAD3), чьи рецепторы и передатчики сигнала попадают в окно \(N \approx 125\)--\(145\).
	\end{itemize}
	
	Это совпадение не случайно. Статус метилирования CpG — это локальная запись координационной истории хроматинового домена. Когда кодируемый белок является узлом универсальной массовой лестницы, его производство и деградация ограничены той же координационной эффективностью \(K_{\mathrm{eff}}\), которая управляет поддержанием эпигенетической метки. Следовательно, возрастной дрейф метилирования сообщает о прогрессирующей потере координации во всей протеомной сети.
	
	\subsection{Фальсифицируемые предсказания}
	\begin{enumerate}
		\item \textbf{Слепой тест на PDB.} Гистограмма \(N_{\mathrm{raw}}\) для всех мономерных белков в Банке данных белков (\(>200\,000\) структур) должна демонстрировать пики на целых и полуцелых значениях с впадинами между ними. Основная биофизика предсказывает гладкое, лишённое особенностей распределение. Этот тест может быть выполнен немедленно с общедоступными данными и не требует новых экспериментов.
		\item \textbf{Гидратационный сдвиг.} Масса белка in vivo увеличивается за счёт первой гидратной оболочки (прочно связанной воды). YUCT предсказывает, что эта дополнительная масса, обычно \(0,3\)--\(0,4\) кДа на кДа белка, сдвигает сырой индекс на величину, сравнимую с топологическим зазором \(\sigma = 0,20\). Это означает, что функциональная координационная глубина белка \emph{определяется} его гидратированной массой, а не массой в вакууме. Высокоточная масс-спектрометрия интактных комплексов в сочетании с измерениями гидродинамического радиуса может проверить это предсказание.
		\item \textbf{Дизайн синтетических белков.} Полипептидная цепь, сконструированная так, чтобы её масса точно соответствовала полуцелому узлу, должна проявлять повышенную термостабильность и более длительный клеточный период полураспада по сравнению с вариантами, смещёнными на \(\pm 0,3\) единицы. Это можно проверить, создав небольшую библиотеку мутантов модельного белка (например, убиквитина или GFP) и измерив скорость их деградации в клетках HeLa.
	\end{enumerate}
	
	Эти предсказания количественны, не содержат свободных параметров и непосредственно фальсифицируемы. Их подтверждение установило бы, что координационная лестница, открытая в физике частиц и ядерной физике, непрерывно продолжается в молекулярные машины жизни.
	
	\begin{figure}[H]
		\centering
		\begin{tikzpicture}
			\begin{axis}[
				width=0.85\textwidth, height=0.55\textwidth,
				xlabel={Полуцелый индекс \(N_f\)},
				ylabel={\(\ln(m/m_e)\)},
				grid=major,
				legend pos=north west,
				legend cell align=left,
				legend style={font=\footnotesize},
				title={Универсальная массовая лестница: от электрона до комплекса ядерной поры},
				title style={font=\bfseries},
				xtick={-10,0,...,200},
				minor xtick={-9.5,-8.5,...,199.5},
				ymin=-2, ymax=30,
				xmin=-5, xmax=195,
				]
				
				% Теоретическая линия: наклон = ln(q)
				\addplot[domain=-10:200, samples=2, dashed, thick, black!50] 
				{x * 0.135155};
				\addlegendentry{Теория: \(\ln m = N_f \ln q\)}
				
				% Фермионы
				\addplot[only marks, mark=*, red, mark size=1.8] coordinates {
					(0, 0)        % e
					(10.5, 1.42)  % u
					(16.5, 2.23)  % d
					(38.5, 5.20)  % s
					(39.5, 5.34)  % mu
					(58.0, 7.84)  % c
					(60.0, 8.11)  % tau
					(66.5, 8.99)  % b
					(94.0,12.71)  % t
				};
				\addlegendentry{Фермионы}
				
				% Лёгкие ядра (для демонстрации непрерывности)
				\addplot[only marks, mark=square*, blue, mark size=1.5] coordinates {
					(55.5, 7.50)  % p
					(58.5, 7.91)  % D
					(64.0, 8.66)  % 4He
					(70.5, 9.53)  % 12C
					(74.5,10.07)  % 16O
				};
				\addlegendentry{Ядра}
				
				% Отобранные белки (главные герои)
				\addplot[only marks, mark=triangle*, green!60!black, mark size=2.2, thick] coordinates {
					(122.5, 16.56) % Убиквитин
					(125.5, 16.97) % Цитохром c
					(126.5, 17.11) % Лизоцим
					(131.5, 17.78) % GFP
					(134.0, 18.12) % GAPDH
					(137.5, 18.59) % Гемоглобин
					(146.0, 19.74) % Нуклеосома
					(147.0, 19.88) % Сплайсосома U1
					(150.5, 20.35) % АТФ-синтаза F1
					(164.0, 22.18) % 70S рибосома
					(164.5, 22.24) % 26S протеасома
					(187.5, 25.36) % NPC
				};
				\addlegendentry{Белковые комплексы}
				
				% Подписи к ключевым биологическим объектам
				\node[green!60!black, font=\tiny, anchor=south west] at (axis cs:164.0,22.18) {Рибосома};
				\node[green!60!black, font=\tiny, anchor=south west] at (axis cs:164.5,22.24) {Протеасома};
				\node[green!60!black, font=\tiny, anchor=south west] at (axis cs:187.5,25.36) {NPC};
				\node[green!60!black, font=\tiny, anchor=south west] at (axis cs:122.5,16.56) {Убиквитин};
				\node[green!60!black, font=\tiny, anchor=south east] at (axis cs:137.5,18.59) {Гемоглобин};
				
				% Вертикальные пунктирные линии на полуцелых значениях (явные команды)
				\draw[gray, thin, dotted] (axis cs:122.5,-2) -- (axis cs:122.5,30);
				\draw[gray, thin, dotted] (axis cs:125.5,-2) -- (axis cs:125.5,30);
				\draw[gray, thin, dotted] (axis cs:126.5,-2) -- (axis cs:126.5,30);
				\draw[gray, thin, dotted] (axis cs:131.5,-2) -- (axis cs:131.5,30);
				\draw[gray, thin, dotted] (axis cs:134,-2) -- (axis cs:134,30);
				\draw[gray, thin, dotted] (axis cs:137.5,-2) -- (axis cs:137.5,30);
				\draw[gray, thin, dotted] (axis cs:146,-2) -- (axis cs:146,30);
				\draw[gray, thin, dotted] (axis cs:147,-2) -- (axis cs:147,30);
				\draw[gray, thin, dotted] (axis cs:150.5,-2) -- (axis cs:150.5,30);
				\draw[gray, thin, dotted] (axis cs:164,-2) -- (axis cs:164,30);
				\draw[gray, thin, dotted] (axis cs:164.5,-2) -- (axis cs:164.5,30);
				\draw[gray, thin, dotted] (axis cs:187.5,-2) -- (axis cs:187.5,30);
			\end{axis}
		\end{tikzpicture}
		\caption{\textbf{Универсальная массовая лестница простирается от фундаментальных частиц до крупнейших клеточных машин.} 
			Каждая точка представляет натуральный логарифм массы объекта относительно электрона. 
			Пунктирная линия — теоретическое предсказание YUCT: массы растут как степени кванта \(q = (3/2)^{1/3} \approx 1,1447\). 
			Красные кружки: элементарные фермионы; синие квадраты: лёгкие атомные ядра (показаны для непрерывности); зелёные треугольники: белковые и рибонуклеопротеидные комплексы, обсуждаемые в тексте. 
			Все точки лежат чрезвычайно близко к прямой, демонстрируя, что один и тот же дискретный масштабный закон управляет как субатомной, так и надмолекулярной материей. 
			Зелёные объекты образуют отчётливый кластер в диапазоне \(N_f \approx 120\)--\(190\), соответствующем окну масс, в котором становятся возможными автономное сворачивание и скоординированная клеточная функция. 
			Этот рисунок даёт структурную основу для часов BCLM‑EC: если все эти машины являются узлами единой координационной лестницы, их функциональное состояние — а значит, и сообщающие о нём метки метилирования — должно следовать одному и тому же универсальному закону ошибки.}
		\label{fig:mass_ladder_bclm}
	\end{figure}
	
	\section{Интерпретация весов}
	\label{sec:weights}
	Положительный вес в таблице \ref{tab:cpgs} означает, что сайт становится гиперметилированным по мере снижения \(\REL\) (т.е. с возрастом); отрицательный вес указывает на гипометилирование. Абсолютная величина отражает чувствительность сайта к координационному состоянию. Примечательно, что сайты с наибольшими весами расположены вблизи генов, вовлечённых в:
	\begin{itemize}
		\item \textbf{Репарацию ДНК и стабильность генома} (например, BRCA1, FANCI, RAD51C через NOTCH1 и TCF25);
		\item \textbf{Митохондриальную функцию и окислительный стресс} (например, TXNL1, CYP2E1);
		\item \textbf{Протеостаз} (например, HSPA13, PDCD4);
		\item \textbf{Клеточную адгезию и сигналинг внеклеточного матрикса} (например, ITGB1, COL18A1).
	\end{itemize}
	Это обогащение согласуется с четырьмя уровнями Протокола долгой жизни (Приложение BCLM), позволяя предположить, что дрейф метилирования, улавливаемый часами, является репортёром того же самого глубинного снижения координации, которому протокол призван противодействовать.
	
	\section{Предсказания и фальсифицируемость}
	\label{sec:predictions}
	
	\begin{enumerate}
		\item \textbf{Замедление старения при обработке по Протоколу долгой жизни.}
		Если Протокол долгой жизни повышает клеточную \(\REL\) в \(f = \REL_{\mathrm{LL}}/\REL_{\mathrm{WT}}\) раз, вероятность ошибки падает в \(f^{-\beta}\) раз. Следовательно, дрейф метилирования на каждом часовом CpG должен замедляться в то же число раз. Для обработки продолжительностью \(t\) лет часы покажут возраст, меньший на
		\[
		\Delta\mathrm{Age} = \frac{1}{\beta\gamma} \ln f.
		\]
		При \(\beta\gamma = 0,20\) (Приложение P) и оценке \(f \approx 1,33\) (основанной на совокупном снижении ошибок в Протоколе долгой жизни) ожидается \(\Delta\mathrm{Age} \approx 1,4\) года после одного месяца обработки. Более длительные обработки приведут к б\'{о}льшим сдвигам. Это предсказание можно проверить, культивируя сконструированные по LL кардиомиоциты параллельно с контролем дикого типа и сравнивая их BCLM‑EC возраст через равные промежутки времени.
		
		\item \textbf{Температурная зависимость скорости часов.}
		Поскольку \(\REL(T) \propto T^{-\gamma}\), клетки, культивируемые при более низкой температуре, имеют более высокую \(\REL\) и, следовательно, более медленный дрейф метилирования. Например, снижение температуры культивирования с \(37^\circ\text{C}\) до \(30^\circ\text{C}\) должно привести к замедлению старения примерно на \(15\%\) за несколько пассажей. Это прямое следствие температурного масштабирования координационной эффективности (Приложения BCLM и P), и его можно проверить в любом типе клеток, для которого откалиброваны BCLM‑EC.
		
		\item \textbf{Дисперсия на уровне единичных клеток.}
		На уровне единичных клеток, в предположении независимости ошибок метилирования между клетками, распределение значений метилирования в заданном часовом CpG должно быть логнормальным с дисперсией, растущей как \(\REL^{-2\beta}\). Секвенирование бисульфитной обработки единичных клеток стареющей популяции может, таким образом, дать независимую оценку \(\REL\), которая должна согласовываться со значением, полученным из массового метилирования.
	\end{enumerate}
	
	\section{Ограничения и открытые вопросы}
	\label{sec:limitations}
	\begin{itemize}
		\item Показатель чувствительности \(S_{\mathrm{BCLM}}\) в настоящее время опирается на возрастной наклон \(\beta_i\) как на прокси для истинной производной \(\partial m/\partial \REL\). Более прямое измерение \(\REL\) в единичных клетках (например, через метаболический поток и активность репарации) обеспечило бы более прочную основу.
		\item Часы были обучены на цельной крови; тканеспецифичные часы должны калиброваться отдельно, так как \(\alpha_{\mathrm{epi}}\) и \(\gamma\) могут различаться.
		\item Модель экспоненциального спада (\ref{eq:Kdecay}) является приближением первого порядка. Более детальные модели деградации координационной сети могут давать неэкспоненциальные формы, которые можно будет различить с помощью продольных данных высокого разрешения.
		\item Веса в таблице \ref{tab:cpgs} до некоторой степени специфичны для платформы Illumina 450K. Будущие версии следует адаптировать для анализов метилирования на основе секвенирования.
	\end{itemize}
	
	\section{Заключение}
	\label{sec:conclusion}
	Координационные часы BCLM предоставляют механистическую, основанную на теории альтернативу чисто статистическим эпигенетическим часам. Они укоренены в универсальном законе ошибки YUCT, отбирают CpG по их чувствительности к клеточному координационному состоянию, достигают конкурентоспособной точности предсказания возраста с небольшим числом сайтов и делают конкретные, фальсифицируемые предсказания о том, как траектории метилирования реагируют на воздействия, изменяющие \(K_{\mathrm{eff}}\). Вместе с Протоколом долгой жизни BCLM‑EC образуют замкнутый экспериментальный цикл: теория предсказывает эпигенетический результат, а результат либо подтверждает, либо опровергает теорию. Это превращает изучение эпигенетического старения из описательной науки в объяснительную.
	
	\vspace{0.5cm}
	\noindent\textbf{Данные и код:} Все данные метилирования взяты из GEO (GSE40279). Обученные коэффициенты часов (таблица \ref{tab:cpgs}), код на R для калибровки и список первоначально рассмотренных 100 CpG доступны по адресу \url{https://github.com/Alexey-Yakushev-YUCT/YPSDC}.
	
	\begin{thebibliography}{99}
		\bibitem{Horvath2013} S. Horvath, ``DNA methylation age of human tissues and cell types,'' \emph{Genome Biology}, 14, R115 (2013).
		\bibitem{Hannum2013} G. Hannum \emph{et al.}, ``Genome-wide methylation profiles reveal quantitative views of human ageing rates,'' \emph{Molecular Cell}, 49, 359--367 (2013).
		\bibitem{Levine2018} M. E. Levine \emph{et al.}, ``An epigenetic biomarker of ageing for lifespan and healthspan,'' \emph{Aging}, 10, 573--591 (2018).
		\bibitem{BCLM} A.\,V.\,Yakushev, ``Appendix BCLM: Biological Coordination Lagrangian,'' in \emph{YUCT}, Zenodo (2026).
		\bibitem{AppendixP} A.\,V.\,Yakushev, ``Appendix P: Temperature Dependence of Gravity,'' in \emph{YUCT}, Zenodo (2026).
		\bibitem{YPSDC_repo} A.\,V.\,Yakushev, YPSDC Repository, \url{https://github.com/Alexey-Yakushev-YUCT/YPSDC}.
	\end{thebibliography}
	
\end{document}